\documentclass[sigconf]{acmart}


%for ACM default CCS concepts
\settopmatter{printfolios=true}%, printacmref=false}
\copyrightyear{2026}
\acmYear{2026}
\acmConference[DS Research Project]{Data Science Research Project}{March 2026}{Sri Lanka}
\acmBooktitle{Data Science Research Project}
\acmDOI{10.1145/1122445.1122456}
\acmISBN{978-1-4503-XXXX-X/26/03}
\PassOptionsToPackage{hyphens}{url}
\usepackage{hyperref}




%\setcopyright{none}

\pagestyle{plain}
%\renewcommand\footnotetextcopyrightpermission[1]{}

\usepackage{comment}
\usepackage{enumitem}
\setlist{noitemsep, topsep=2pt}
\sloppy
\usepackage{titlesec}
\titlespacing*{\section}{0pt}{4pt}{2pt}
\titlespacing*{\subsection}{0pt}{3pt}{1pt}
\titlespacing*{\subsubsection}{0pt}{2pt}{1pt}

\PassOptionsToPackage{hyphens}{url}
\usepackage{hyperref}
\usepackage{booktabs}
\usepackage{multirow}
\usepackage{caption}

% Tighten space around floats (tables)
\setlength{\intextsep}{4pt plus 1pt minus 1pt}
\setlength{\abovecaptionskip}{3pt}
\setlength{\belowcaptionskip}{1pt}
\setlength{\textfloatsep}{6pt plus 1pt minus 2pt}

\usepackage{multicol} %references section two columns

\begin{document}

\title{A Data-Driven Analysis of Migration and Remittance: Insights for Related Decision-Making in Sri Lanka}

\author{Ginige D.N.}
\email{dinurag.23@cse.mrt.ac.lk}
\affiliation{\institution{University of Moratuwa}\streetaddress{Dept. of Computer Science \& Engineering}\city{Moratuwa}\country{Sri Lanka}}

\author{Lakshan H.M.K.}
\email{kalanal.23@cse.mrt.ac.lk}
\affiliation{\institution{University of Moratuwa}\streetaddress{Dept. of Computer Science \& Engineering}\city{Moratuwa}\country{Sri Lanka}}

\author{Fernando K.N.D.}
\email{dhinanjayaf.23@cse.mrt.ac.lk}
\affiliation{\institution{University of Moratuwa}\streetaddress{Dept. of Computer Science \& Engineering}\city{Moratuwa}\country{Sri Lanka}}

\author{Gurusinghe C.R.}
\email{chanupag.23@cse.mrt.ac.lk}
\affiliation{\institution{University of Moratuwa}\streetaddress{Dept. of Computer Science \& Engineering}\city{Moratuwa}\country{Sri Lanka}}

\author{Pahanga G.H.L.}
\email{lasanap.23@cse.mrt.ac.lk}
\affiliation{\institution{University of Moratuwa}\streetaddress{Dept. of Computer Science \& Engineering}\city{Moratuwa}\country{Sri Lanka}}

\begin{abstract}
This study applies data science and machine learning to 1994--2025 Sri Lankan migration and remittance data across three integrated analytical paths: demographic, macroeconomic, and time series. Using 384 observations(32yrs x 12months) from eight authoritative sources, exploratory, econometric (STL, Johansen cointegration, VAR(12), VECM, Granger causality, Chow tests), and predictive (ARIMA, Ridge, Gradient Boosting) methods identify key structural and macroeconomic drivers of remittance dynamics. USD/LKR, CCPI, and oil price are confirmed as significant Granger causes of remittance variation; two cointegrating equilibria and three structural breaks (2008, 2020, 2022) are formally verified. An evidence-based 2026 forecast of USD~8,148~Mn annually (MAPE~7.7\%) is produced, with findings informing exchange-rate policy, migrant welfare, and sustainable development planning.
\end{abstract}

\keywords{Migration, Remittances, Data Science, Machine Learning, Time Series Analysis, Granger Causality, Sri Lanka}

\maketitle

\section{Introduction}
%==================================================================================
Remittances constitute one of the most significant and stable sources of external finance for developing economies. For Sri Lanka, worker remittances have consistently accounted for 6--9\% of GDP~\cite{remittance-gdp}, making them a dominant source of foreign exchange. Beyond macroeconomic stabilisation, remittances fund household consumption, education, and healthcare persisting through crises, including the 2022 sovereign debt default when the government relied on diaspora transfers to service essential imports~\cite{ilo-migration}.

Existing research on Sri Lankan migration and remittances is largely descriptive or relies on narrow, cross-sectional econometric methods with short observation windows, annual data, and a failure to address non-stationarity or integrate demographic and macroeconomic dimensions simultaneously~\cite{ramanayake2018,ilo-migration}. This study addresses all four gaps. Specifically, we contribute: (i)~a 32-year dataset of 384 observations interpolated to monthly frequency from eight authoritative sources; (ii)~stationarity-corrected causal inference via Granger tests and VECM that distinguishes genuine from spurious predictors; (iii)~formal structural break detection at three historically motivated dates; (iv)~an integrated three-path analytical framework connecting demographic composition, macroeconomic conditions, and temporal patterns; and (v)~a 2026 annual remittance forecast of USD~8,148~Mn (MAPE~7.7\%), outperforming ARIMA by a factor of 3.7.

%==================================================================================
\section{Data Sources and Methodology}
%==================================================================================

\subsection{Data Sources and Collection}
Data spanning 1994--2025 were sourced from eight authoritative institutions: \textbf{Central Bank of Sri Lanka (CBSL)}~\cite{cbsl} (remittances, exchange rates, financial statistics); \textbf{World Bank Open Data}~\cite{worldbank-remittances,worldbank-migration,worldbank-ppp} (development indicators, GDP PPP); \textbf{SLBFE}~\cite{slbfe} (monthly outward labour migration by skill and gender); \textbf{ILO}~\cite{ilo} (labour statistics); \textbf{FRED}~\cite{fred-remittances,cbsl-exchangerate} (exchange rates); \textbf{DCS Sri Lanka}~\cite{ccpi} (CCPI, population); \textbf{U.S.~EIA}~\cite{eia-oil} (oil prices); and \textbf{EM-DAT}\cite{emdat2026} (disaster counts). Collection employed automated web extraction, OCR-based PDF digitisation of SLBFE annual reports (1994--2005), and structured CSV/XLSX imports. All indicators were validated against at least one independent source before inclusion.

\subsection{Analytical Framework}
Three interconnected analytical paths form the methodological backbone:
\begin{enumerate}[noitemsep,topsep=0pt]
\item \textbf{Demographic Path:} Gender composition transitions, skill polarisation, district-level concentration, age--contract correlations, and poverty--emigration linkages using 32 annual observations.
\item \textbf{Macroeconomic Path:} GDP PPP, CCPI, LKR/USD exchange rate, unemployment, oil prices, disaster impact, and transfer costs as drivers of both emigration and remittance behaviour.
\item \textbf{Time Series Path:} STL decomposition, ADF/KPSS stationarity testing, Johansen cointegration, VAR/VECM, Granger causality, impulse response functions, FEVD, Chow structural breaks, and ML forecasting on 384 observations.
\end{enumerate}
Phase~2 integration synthesised cross-path findings through rolling correlations, cross-correlation functions, and a unified predictive modelling pipeline. Each path was first executed independently before cross-path variables were introduced, ensuring individual analytical conclusions remained valid prior to integration.

\subsection{Datasets}
Three primary datasets were constructed (Table~\ref{tab:datasets-summary}); Dataset~3 (384 observations at monthly frequency) forms the core of the econometric and ML analyses. 

\begin{minipage}{\columnwidth} % table and table 1 name binding
{\small
\captionof{table}{Summary of Datasets and Variables}
\label{tab:datasets-summary}
\begin{tabular}{|p{1.8cm}|p{1.2cm}|p{4.6cm}|}
\hline
\textbf{Dataset} & \textbf{Period} & \textbf{Attributes} \\
\hline
Demographic (DS1) & 1994--2025 & Total Emigration, Gender Split, Skill Level, Region, Avg.\ Age, Contract Duration \\
\hline
Macroeconomic (DS2) & 1994--2025 & GDP PPP, CCPI, LKR/USD, Unemployment, Oil Price, Disaster Count, Transfer Cost \\
\hline
Time Series (DS3) & 1994--2025 & Remittances (USD Mn/yr), SLBFE Departures, LKR/USD, Oil Price, CCPI, GDP PPP, Skilled\% \\
\hline
\end{tabular}
}
\end{minipage} % table and table 1 name binding
\vspace{-6pt}

\subsection{Data Cleaning and Preprocessing}

\subsubsection{Cleaning Pipeline}
Missing values were handled hierarchically: columns exceeding 40\% missingness were dropped or replaced with proxy variables; isolated gaps ($<$5\%) received median/mean imputation via \texttt{SimpleImputer}; consecutive annual gaps (prevalent in 1994--2000 SLBFE records) were filled by linear interpolation; and temporal series gaps were addressed with direction-conditioned forward/backward fill. Outliers were flagged using Z-score ($|z|>3$) and confirmed by visual inspection, then cross-referenced against historical events: crisis-year observations were retained as analytically significant while data-entry errors were corrected via interpolation or source cross-validation. Categorical inconsistencies were standardised; logical constraints (male\%$+$female\%$=$100) were enforced; and duplicate records were removed with \texttt{drop\_duplicates()}.

\subsubsection{Harmonisation, Normalisation and Feature Engineering}
Remittance data originate as World Bank annual totals (USD Mn/year) and were linearly interpolated to monthly frequency using $m(t) = y_n + [y_{n+1}-y_n]\times\frac{t-t_n}{t_{n+1}-t_n}$, with each year's anchor placed at January. All other variables were harmonised analogously: exchange rates via within-month averaging; SLBFE departures via uniform monthly allocation; and CCPI via chain-linking. All variables were Z-score normalised ($\mu=0,\,\sigma=1$); right-skewed variables (GDP PPP, wages) were log-transformed prior to normalisation. Stationarity was confirmed via ADF and KPSS tests; all series were $I(1)$, and first-differencing ($\Delta Y_t=Y_t-Y_{t-1}$) was applied for all correlation, Granger, and VAR analyses. A binary classification target was constructed by thresholding annual remittances at the sample median (USD~3,978~Mn). STL decomposition (period~$=$~12, robust) extracted trend and residual components from the annual remittance series; no meaningful seasonal component is present given the annual data source. For ML, the feature matrix was augmented with lagged first-differenced remittance changes at 1, 3, 6, and 12 months, a composite Gulf pull-factor, and skill-composition ratios. Pearson screening ($|r|>0.9$) found no redundant pairs; three variables with VIF~$>$10 (CCPI, GDP, wages) were retained for their distinct analytical roles.

%==================================================================================
\section{Analysis and Results}
%==================================================================================

\subsection{Exploratory Data Analysis}

\noindent\textbf{Gender Composition Transition:}
From 1994 to 2025, outward labour migration shifted from a predominantly female flow to a 65:35 male:female ratio by the late 2010s. Male participation increased markedly post-2008 as Gulf construction and technical roles expanded. Despite this aggregate shift, internal skill composition remained highly gendered: female emigrants stayed concentrated in low-skilled categories ($\sim$87\% of all female emigrants in 2017 were low-skilled), while males maintained a more balanced skilled/unskilled split. A linear regression on the annual female emigrant share confirmed a statistically significant secular decline.
%$($\beta=-1.271$, $p<0.001$, $R^2=0.84$).

\noindent\textbf{Skill Polarisation and Human Capital Drift:}
Overall skilled emigrant participation declined at $\sim$0.44~pp/year over the full sample - female skilled share falling ($-0.25$~pp/year), male rising ($+0.49$~pp/year). Structural breaks were detected in 2018 and 2020, followed by a strong post-2020 re-skilling trend ($+5.19$~pp/year). A one-sided Wilcoxon signed-rank test across 32 annual year-pairs (Shapiro--Wilk non-normality: $p=0.013$) found a substantial and persistent gender disparity in skill level composition. 

\noindent\textbf{Poverty--Migration Link:}
A non-linear, temporally lagged relationship links poverty and emigration. During 1994--2000, high poverty rates ($>$60\%) coincided with elevated emigration consistent with a survival-migration motive. As poverty fell to 18.6\% by 2017, emigration fluctuated rather than receded indicating migration hysteresis: rising household income expands access to legal migration channels and raises aspirations even as material need recedes, with a 1--2 year lag between economic shifts and migration responses. A Pearson test across 32 annual observations yielded $r=-0.328$ ($t=-1.90$, $p=0.067$, $R^2=0.107$), failing to reject the null at 5\% but significant directionally at 10\%.

\noindent\textbf{District Concentration and Mobility:}
Regional analysis reveals intensifying geographic concentration: normalised HHI rises $\sim$9~points/year with strong rank persistence (Markov diagonal probabilities 0.84--1.00) and no detected structural break, implying historically dominant emigration districts continuously reinforce their lead.

\noindent\textbf{Macro--Migration--Remittance Dynamics:}
Among macroeconomic variables in levels, remittances correlate most strongly with GDP PPP ($r=0.946$) and negatively with transfer cost ($r=-0.549$); migration departures are most linked to unemployment ($r=-0.704$). Disaster months differ significantly from non-disaster months in migration ($t=2.55$, $p=0.012$) and remittances ($t=3.68$, $p<0.001$). A Pearson correlation between annual SLBFE departures and remittance inflows confirmed a moderate association ($r=0.600$, $p<0.001$, $R^2=0.360$). Older emigrants also secure longer contracts ($r=0.591$, $p<0.001$, 95\% CI~[0.303, 0.777]).

\subsection{Time Series Econometric Analysis}

\noindent\textbf{Descriptive Statistics and STL Decomposition:}
The annual remittance series (interpolated to 384 observations) exhibits strong secular growth: sample mean USD~4,046~Mn/year, trough USD~715~Mn (1994), peak USD~8,069~Mn (2025). STL decomposition (period~$=$~12, robust fitting) identifies: (i)~a strong upward trend (USD~718~Mn in 1994 to USD~8,256~Mn in 2025); and (ii)~pronounced residual spikes at crisis dates, with the 2022 economic crisis generating the largest negative irregular component ($\approx-951$~USD~Mn). No seasonal component is interpreted as the data originates from annual WB figures.

\noindent\textbf{Stationarity (ADF/KPSS):}
ADF and KPSS tests confirmed all six key series as $I(1)$ non-stationary in levels, stationary in first differences. PACF of first-differenced remittances shows significant spikes at lags 1--2, motivating ARIMA$(2,1,0)$. All inferential analyses use first-differenced series exclusively to avoid spurious inference.

\noindent\textbf{Johansen Cointegration:}
The Johansen trace test on the five-variable system \{remittances, USD/LKR, oil price, CCPI, GDP PPP\} identified $r=2$ cointegrating vectors at the 95\% level (Table~\ref{tab:johansen}), confirming long-run equilibrium relationships and mandating VECM.

\begin{table}[h]
\centering
\caption{Johansen Trace Test Results}
\label{tab:johansen}
\small
\begin{tabular}{lrrr}
\toprule
\textbf{H\textsubscript{0}} & \textbf{Trace Stat} & \textbf{CV 95\%} & \textbf{Integrated?} \\
\midrule
$r \leq 0$ & 222.36 & 69.82 & Yes \\
$r \leq 1$ & 65.18  & 47.85 & Yes \\
$r \leq 2$ & 22.72  & 29.80 & No  \\
$r \leq 3$ & 3.89   & 15.49 & No  \\
\bottomrule
\multicolumn{4}{l}{\small 2 cointegrating vectors confirmed $\Rightarrow$ VECM preferred over VAR.}
\end{tabular}
\end{table}

\noindent\textbf{Cross-Correlation and Granger Causality:}
A VAR(12) model (lag order selected by AIC, AIC$_{\min}$ at lag~12) was estimated on first-differenced series. Granger causality F-tests confirm USD/LKR ($p=0.0001$), CCPI ($p=0.008$), and oil price ($p=0.010$) as statistically significant short-run predictors of remittance changes; GDP PPP was not significant ($p=0.208$). Cross-correlation on first-differenced series reveals meaningful lead-lag structure: CCPI leads remittances by 14 observations ($r=-0.377$), USD/LKR leads by 5 observations ($r=-0.322$), oil price leads by 5 observations ($r=-0.163$), while GDP PPP lags remittances ($r=+0.582$, lag~$+24$) suggesting per-capita income growth responds to rather than drives remittance flows.

\noindent\textbf{VECM Short-Run Dynamics:}
The VECM on $r=2$ vectors reveals significant short-run predictors of remittance changes: the lagged remittance term at lag~1 ($\hat{\beta}=0.807$, $p<0.001$) dominates, confirming strong path dependence; USD/LKR at lag~1 ($\hat{\beta}=-1.079$, $p=0.025$) carries a significant negative effect; and oil price at lag~2 ($\hat{\beta}=-0.745$, $p=0.001$) shows a delayed negative impact on remittance inflows.

\noindent\textbf{Structural Breaks (Chow Test):}
All three candidate break dates were formally confirmed at $p<0.001$ (Table~\ref{tab:chow}). The COVID-19 onset produced the largest disruption (F~$=$~156.5), reflecting simultaneous collapse of departure volumes and remittance channels.

\begin{table}[h]
\centering
\caption{Chow Structural Break Tests on Remittance Series}
\label{tab:chow}
\small
\begin{tabular}{lrr}
\toprule
\textbf{Break Event} & \textbf{F-Statistic} & \textbf{p-value} \\
\midrule
2008 Financial Crisis      & 130.35 & $<0.001$ \\
COVID-19 Onset (Mar 2020)  & 156.47 & $<0.001$ \\
Economic Crisis (Apr 2022) &  46.59 & $<0.001$ \\
\bottomrule
\end{tabular}
\end{table}

\noindent\textbf{Trend-Adjusted Shock Analysis:}
STL trend-adjusted 12-month pre/post windows isolate the shock-attributable deviation from the secular trend. The 2008 GFC raised remittances by +18.0\% (diaspora solidarity and precautionary-transfer response). COVID-19 reduced them by $-11.4\%$ (supply-side disruption to labour mobility). The 2022 economic crisis produced a counter-cyclical +8.4\% increase, consistent with informal channel substitution and distress-motivated household transfers as the foreign-exchange shortage intensified.

\subsection{Predictive Analysis}

\subsubsection{ML Task Formulation}
Two complementary tasks were implemented on temporally ordered, stationarity-corrected data: (1)~\textbf{Regression} predicting annual remittance changes from lagged autoregressive features and stationary macro predictors; (2)~\textbf{Binary Classification} labelling each observation as high ($>$~sample median USD~3,978~Mn/yr) or low remittance using the same feature matrix. 

\subsubsection{Evaluation Metrics and Validation Strategy}
Supervised regression models were evaluated under five-fold expanding-window time-series cross-validation preserving temporal order and preventing look-ahead bias with RMSE, MAE, and $R^2$ as primary metrics. Classification performance was assessed by accuracy and weighted F1-score.

\subsubsection{Machine Learning Results}

\noindent\textbf{Supervised Regression:}
Ridge regression substantially outperformed Lasso and Gradient Boosting on the stationary lagged feature matrix (Table~\ref{tab:ml-results}). Feature importance confirms overwhelming autoregressive dominance: $\Delta$remittance lag-1 alone accounts for $\approx$78\% of predictive signal, with lag-12 contributing $\approx$17\%. GDP PPP change contributes $\approx$4\%, while USD/LKR and SLBFE departures together contribute less than 1\% confirming that near-term remittance changes are primarily driven by their own recent history rather than contemporaneous macro shifts.

\begin{table}[h]
\centering
\caption{Supervised Model Performance (5-fold Time-Series CV)}
\label{tab:ml-results}
\small
\begin{tabular}{lccc}
\toprule
\textbf{Model} & \textbf{RMSE} & \textbf{MAE} & \textbf{$R^2$} \\
\midrule
Ridge Regression  & 23.46 & 12.12 & \phantom{$-$}0.532 \\
Lasso Regression  & 24.37 & 12.88 & \phantom{$-$}0.217 \\
Gradient Boosting & 45.58 & 33.80 & $-$0.521 \\
\bottomrule
\end{tabular}
\end{table}

\noindent\textbf{Binary Classification:}
A Decision Tree classifier achieved 67.9\% accuracy and weighted F1-score of 0.750 distinguishing high vs.\ low remittance years under time-series cross-validation, confirming that stationary lagged features carry directional predictive signal above a naive 50\% baseline.

\subsubsection{Forecasting Results}

\noindent\textbf{ARIMA(2,1,0):} Fitted on 1994--2019 (312 observations), AR order from PACF analysis. The fitted AR$_1=0.886$ ($p<0.001$) and AR$_2=0.094$ ($p=0.318$). Out-of-sample performance on the 2020--2025 test set was poor (RMSE~$=$~1,964, MAPE~$=$~31.0\%), with ARIMA unable to adapt to three structural breaks within the evaluation window. ARIMA's poor performance is consistent with the Chow test results confirming breaks at 2020 and 2022, the model's stationarity assumption is violated precisely during the evaluation period.

\noindent\textbf{ML Forecasting:} The lagged regression model substantially outperformed ARIMA on the same 2020--2025 test set: RMSE~$=$~525, MAPE~$=$~7.7\%. Gradient Boosting achieved RMSE~$=$~819, MAPE~$=$~11.7\%. The lagged regression's strong performance reflects the dominant autoregressive structure: lag-1 alone explains $\sim$78\% of annual remittance variation.

\noindent\textbf{2026 Forecast:} Re-estimated on the full 384-observation sample, the preferred lagged regression projects a 2026 annual total of \textbf{USD~8,148~Mn} (range: 8,079--8,226~Mn; 95\% CI: $\pm$146~Mn), representing \textbf{+1.6\% growth} over the 2025 figure of USD~8,069~Mn, a modest continuation of the recovery trajectory consistent with the stabilising post-crisis growth rate.

\subsubsection{Statistical Inference Summary}
All five formal hypotheses are summarised in Table~\ref{tab:hypotheses}. Four of five nulls were rejected, confirming a gender composition transition, a persistent gender--skill gap, an age--contract association, and a positive departure--remittance link. Poverty as a direct push factor (H4) was not confirmed at 5\%, supporting a capabilities-based migration framework.

\begin{table}[h]
\centering
\caption{Summary of Hypothesis Tests}
\label{tab:hypotheses}
\small
\begin{tabular}{p{0.3cm}p{3.4cm}p{1.3cm}p{1.3cm}p{0.7cm}}
\toprule
\textbf{H} & \textbf{Hypothesis} & \textbf{Test} & \textbf{Key Stat} & \textbf{Result} \\
\midrule
1 & Female share declining over time   & Lin.\ Reg. & $\beta{=}{-}1.27$, $p{<}.001$ & Reject \\
2 & Female emigrants have higher skilled share & Wilcoxon & $W{=}528$, $p{<}.001$ & Reject \\
3 & Older emigrants secure longer contracts & Pearson & $r{=}.591$, $p{<}.001$ & Reject \\
4 & Poverty drives emigration volume  & Pearson & $r{=}{-}.328$, $p{=}.067$ & Fail \\
5 & Departures predict remittances     & Pearson & $r{=}.600$, $p{<}.001$ & Reject \\
\bottomrule
\multicolumn{5}{l}{\small H4: null not rejected at 5\%; directionally significant at 10\%.}
\end{tabular}
\end{table}

%==================================================================================
\section{Discussion}
%==================================================================================

Granger causality tests on first-differenced series confirm USD/LKR ($p=0.0001$), CCPI ($p=0.008$), and oil price ($p=0.010$) as genuine short-run predictors. GDP PPP is not a significant Granger cause ($p=0.208$) despite its high levels correlation ($r=0.946$) collapsing to $r=0.019$ in first differences and contributing only $\approx$4\% to ML feature importance, confirming it as a long-run co-trending variable rather than a short-run driver.

The two Johansen cointegrating vectors establish stable long-run equilibria. CCPI leads remittances by 14 observations ($r=-0.377$) and USD/LKR by 5 ($r=-0.322$), providing policy-actionable early-warning windows. GDP PPP lagging remittances (lag~$+24$, $r=+0.582$) suggests reverse causality: remittance inflows support consumption and investment, subsequently boosting per-capita income.

The three confirmed structural breaks (Chow $p<0.001$ at all dates) show divergent shock responses. The 2008 GFC (+18.0\%) and 2022 crisis (+8.4\%) each produced counter-cyclical surges, while the 2020 COVID response ($-11.4\%$) is structurally distinct - a supply-side disruption to migration itself.

The demographic path adds distributional nuance. The confirmed gender composition transition (H1) reflects structural shifts in GCC labour demand from domestic-service to construction and technical roles. The persistent gender--skill disparity (H2: 30.77~pp gap, $p<0.001$) reveals that female emigrants remain disproportionately concentrated in lower-skilled, lower-wage occupations, a structural vulnerability with direct implications for bilateral labour agreements and migrant welfare protections. The failure to reject H4 aligns with a capabilities-based migration theory: rising household incomes simultaneously raise migration aspirations and channel access, sustaining emigration flows even as poverty recedes.

Finally, the forecasting comparison exposes a fundamental limitation of linear ARIMA in structurally unstable environments (MAPE 31.0\%) and demonstrates that the lag-enriched Ridge regression achieves dramatically superior out-of-sample performance (MAPE 7.7\%). The modest +1.6\% growth projection to USD~8,148~Mn in 2026 reflects a stabilising post-crisis trajectory, with the dominant autoregressive lag-1 signal ($\approx$78\% of predictive variance) indicating that the recovery pace is now self-reinforcing rather than shock-driven. The negative $R^2$ of Gradient Boosting ($-$0.521) indicates overfitting on the limited annual sample as tree ensembles require more data than linear models to generalise.

%==================================================================================
\section{Conclusion}
%==================================================================================

This study delivers a multi-path analysis of Sri Lankan migration and remittance dynamics (1994--2025, 384 observations). Key contributions: (1)~USD/LKR, CCPI, and oil price confirmed as genuine short-run Granger causes - GDP PPP shown to be a spurious levels predictor; (2)~two long-run equilibria via Johansen cointegration; (3)~Chow structural breaks at 2008, 2020, and 2022 ($p<0.001$); (4)~counter-cyclical shock responses (+18.0\% in 2008, +8.4\% in 2022) versus COVID disruption ($-11.4\%$); (5)~persistent gender--skill gap (30.77~pp, $p<0.001$); and (6)~2026 forecast of USD~8,148~Mn (MAPE 7.7\%), 3.7 times more accurate than ARIMA.

For policymakers, the findings indicate that exchange rate stabilisation and transfer cost reduction are the highest-leverage interventions for sustaining remittance inflows, while migrant protection frameworks should prioritise female low-skilled workers who remain structurally disadvantaged. Central Bank reserve planners can exploit the strong autoregressive forecast signal and the identified 5--14 observation macro lead-lag structure for forward-looking reserve management. 

Limitations include the use of World Bank annual remittance totals interpolated to monthly frequency, imputation of gender-disaggregated departure data prior to SLBFE digital records, and the absence of bilateral destination-country data. As GDP PPP is also sourced annually, its Granger non-significance ($p=0.208$) should be interpreted as a data-structure sensitivity rather than a definitive economic conclusion, both variables share the same annual-boundary signal concentration. 

\section{Future Work}
Future work will incorporate destination-country labour demand indicators, bilateral migration agreement variables, and SLBFE bilateral records to improve causal identification. The primary priority for the final project phase is a unified VECM-augmented ML forecasting system integrating demographic, macroeconomic, and time series lags from all three analytical paths.



\clearpage
\onecolumn
\begin{multicols*}{2}
\raggedright


\begin{thebibliography}{19}

\bibitem{remittance-gdp}
TheGlobalEconomy.com. 2026.
\textit{Sri Lanka: Remittances Percent of GDP}.
\url{https://www.theglobaleconomy.com/Sri-Lanka/remittances_percent_GDP/}

\bibitem{ilo-migration}
P. Wickramasekara. 2018.
\textit{Sri Lanka's Labour Migration Trends, Remittances and Economic Growth}.
International Labour Organization.
\url{https://www.researchgate.net/publication/327162787}

\bibitem{worldbank-remittances}
World Bank. 2026.
\textit{Personal Remittances, Received: Sri Lanka}.
\url{https://data.worldbank.org/indicator/BX.TRF.PWKR.CD.DT?locations=LK}

\bibitem{worldbank-migration}
World Bank. 2026.
\textit{Net Migration: Sri Lanka}.
\url{https://data.worldbank.org/indicator/SM.POP.NETM?locations=LK}

\bibitem{fred-remittances}
Federal Reserve Bank of St. Louis. 2026.
\textit{Remittance Inflows to GDP for Sri Lanka}.
\url{https://alfred.stlouisfed.org/series?seid=DDOI11LKA156NWDB}

\bibitem{ccpi}
Department of Census and Statistics. 2026.
\textit{Monthly Colombo Consumer Price Index (CCPI)}.
\url{https://www.statistics.gov.lk/InflationAndPrices/StaticalInformation/MonthlyCCPI}

\bibitem{worldbank-ppp}
World Bank. 2026.
\textit{GDP per Capita, PPP: Sri Lanka}.
\url{https://data.worldbank.org/indicator/NY.GDP.PCAP.PP.CD?locations=LK}

\bibitem{cbsl-exchangerate}
Federal Reserve Bank of St. Louis. 2026.
\textit{Sri Lanka Rupee to U.S. Dollar Spot Exchange Rate}.
\url{https://fred.stlouisfed.org/series/DEXSLUS}

\bibitem{eia-oil}
U.S. Energy Information Administration. 2026.
\textit{Brent Crude Oil Prices}.
\url{https://www.eia.gov/dnav/pet/hist/LeafHandler.ashx?n=PET\&s=RBRTE\&f=M}

\bibitem{slbfe}
Sri Lanka Bureau of Foreign Employment. 2026.
\textit{Annual Statistical Handbooks on Outward Labour Migration}.
\url{https://www.slbfe.lk/statistics/}

\bibitem{ramanayake2018}
S. S. Ramanayake and C. S. Wijetunga. 2018.
Sri Lanka's labour migration trends, remittances and economic growth.
\textit{South Asia Research} 38, 4 (2018), 615--618.
\url{https://doi.org/10.1177/0262728018792088}

\bibitem{cbsl}
Central Bank of Sri Lanka. 2026.
Central Bank eResearch Portal.
\url{https://www.cbsl.lk/eResearch/}

\bibitem{ilo}
International Labour Organization. 2018.
Sri Lanka Labour Migration Trends, Remittances and Economic Growth.
\url{https://www.ilo.org/media/334001/download}

\bibitem{johansen1988}
S. Johansen. 1988.
Statistical analysis of cointegration vectors.
\textit{Journal of Economic Dynamics and Control} 12 (1988), 231--254.

\bibitem{engle1987}
R. F. Engle and C. W. J. Granger. 1987.
Co-integration and error correction: Representation, estimation, and testing.
\textit{Econometrica} 55, 2 (1987), 251--276.

\bibitem{granger1969}
C. W. J. Granger. 1969.
Investigating causal relations by econometric models and cross-spectral methods.
\textit{Econometrica} 37, 3 (1969), 424--438.

\bibitem{cleveland1990}
R. B. Cleveland, W. S. Cleveland, J. E. McRae, and I. Terpenning. 1990.
STL: A seasonal-trend decomposition procedure based on loess.
\textit{Journal of Official Statistics} 6, 1 (1990), 3--73.

\bibitem{chow1960}
G. C. Chow. 1960.
Tests of equality between sets of coefficients in two linear regressions.
\textit{Econometrica} 28, 3 (1960), 591--605.

\bibitem{emdat2026}
Center for Research on the Epidemiology of Disasters (CRED). 2026. \textit{EM-DAT: The International Disaster Database}. UCLouvain, Brussels, Belgium. Retrieved February 22, 2026 from https://www.emdat.be

\end{thebibliography}
\end{multicols*}

\end{document}